\documentclass[12pt,a4paper]{ctexart}

\usepackage{ctex}

\usepackage[a4paper,margin=1in]{geometry}

\usepackage{amsmath,amssymb,amsfonts}

\usepackage{graphicx}

\usepackage{setspace}
\onehalfspacing

\usepackage{enumitem}
\setlist{itemsep=2pt, topsep=2pt}

\usepackage{titlesec}

\ctexset{
	section = {
	  format = \Large\bfseries\raggedright
	 },
	subsection = {
			format = \large\bfseries\raggedright
		}
}

\usepackage{hyperref}

\usepackage{cite}


\setcounter{tocdepth}{2}

\begin{document}
\section{研究背景与目前系统现状}
\subsection{研究背景}

对于蛋白质设计这类高代价任务而言，系统真正面对的难点不只是”是否会调用工具“，而在于以下几个方面：

\begin{itemize}
	\item 任务链条长，且阶段之间存在强依赖关系；
	\item 不同阶段的计算代价和失败后果高度不对称；
	\item 运行过程中常常出现信息不充分、候选分歧、局部失败和恢复决策问题；
	\item 何时继续自动执行、何时切换策略、何时请求人工介入，本身就是关键问题。
\end{itemize}

因此，如果将毕设继续定义为“蛋白质设计助手”，其研究重点很容易停留在“任务分解、工具调用、结果输出”这一层面，即便加入了风险门控等系统。如果希望体现更明确的计算机相关的专业特征的话，我觉得可以把问题进一步上升为“高代价科学工作流的执行控制与治理”。

\subsection{当前系统现状}

从现有的系统设计来看，它已经形成了一个较为稳定的多Agent工作流执行系统。

首先，系统采用了显式有限状态机（FSM）作为运行核心，要求任务状态显式、状态转移合法且可记录，外部语义状态和内部执行状态均有明确定义。这意味着系统已经具备较好的执行一致性、可恢复性与可审计性基础。

其次，系统在角色划分上已经形成了较为清晰的职责边界：PlannerAgent负责生成计划候选，ExecutorAgent负责工具执行和恢复流程，SafetyAgent负责评估，SummarizerAgent负责结果汇总。这种结构表明系统已经具备典型多Agent runtime的雏形。

再次，系统已经围绕蛋白质设计任务构建了较为完整的阶段化流程，已经支持：
\begin{itemize}
	\item S1 序列候选生成
	\item S2 结构投影
	\item S3 质量门控
	\item S4 基于结构条件的细化
	\item 分层 patch/replan 恢复机制
\end{itemize}

此外，系统已经具备 HITL（Human-in-the-loop）机制和恢复链设计，包括 \texttt{WAITING\_*} 以及 $retry \to patch \to replan$ 的恢复顺序
\section{当前选题面临的核心问题}
\section{选题目标的重新定义}
\section{候选方向比较}
\section{选择Belief-State Controller}
\section{Belief-State Controller 问题定义}
\section{相关研究现状}
\section{可能的创新点、难点}
\section{可行性分析}
\section{预期结果}
\section{风险与备选方案}
\end{document}
